\documentclass{beamer}

\usepackage{presentation}

\usepackage{minted}
\setminted{
  fontsize=\scriptsize,
  breaklines=true,
}

\title{Elsőrendű logika, mint algebrai elmélet formalizálása Agdában}
\author{Szilágyi Péter \\[1ex] \small Témavezető: Kaposi Ambrus}
\institute{ELTE IK – Programtervező informatikus MSc}
\date{2025}

\begin{document}

\begin{frame}
  \titlepage
\end{frame}

\begin{frame}{A dolgozat célja és motiváció}
  Az elsőrendű logika
  \begin{itemize}
    \item formalizálása Agdában
    \item algebrai elméletként megadva
        \begin{itemize}
        \item algebrai elmélet = szortok, operátorok, egyenletek
        \item szintaxis = iniciális modell
            \begin{itemize}
            \item kvóciensek nélkül megadható
            \end{itemize}
        \end{itemize}
    \item több példa szemantikával
        \begin{itemize}
            \item Bool modell: igaz / hamis
            \item Tarski modell: standard szemantika
            \item Family modell: típuscsalád alapú szemantika
        \end{itemize}
  \end{itemize}
\end{frame}

\begin{frame}[fragile]{Modellfogalom}
    \begin{minted}{agda}
Con     : Set
Sub     : Con → Con → Set
Tm      : Con → Set
For     : Con → Set
Pf      : (Γ : Con) → For Γ → Prop
_[_]ᵗ   : Tm Γ → Sub Δ Γ → Tm Δ
_[_]ᶠ   : For Γ → Sub Δ Γ → For Δ
_[_]ᵖ   : Pf Γ A → (γ : Sub Δ Γ) → Pf Δ (A [ γ ]ᶠ)
_▹ₜ     : Con → Con
_▹ₚ_    : (Γ : Con) → For Γ → Con
...
_⊃_     : For Γ → For Γ → For Γ
⊃[]     : (A ⊃ B) [ γ ]ᶠ = (A [ γ ]ᶠ) ⊃ (B [ γ ]ᶠ)
⊃in     : Pf (Γ ▹ₚ A) (B [ pₚ ]ᶠ) → Pf Γ (A ⊃ B)
⊃out    : Pf Γ (A ⊃ B) → Pf Γ A → Pf Γ B
...
∀       : For (Γ ▹ₜ) → For Γ
∀[]     : (∀ A) [ γ ]ᶠ = ∀ (A [ (γ ∘ pₜ) ,ₜ qₜ ]ᶠ)
∀in     : Pf (Γ ▹ₜ) A → Pf Γ (∀ A)
∀out    : Pf Γ (∀ A) → Pf (Γ ▹ₜ ) A
    \end{minted}
\end{frame}

\begin{frame}[fragile]{Szintaxis}
    \begin{minted}{agda}
data Cont : Set
  _▹ₜ     : Cont → Cont
data Tm   : Cont → Set
  var     : VTm Γ → Tm Γ
  fun     : funar n → Tm Γ ^ n → Tm Γ
data Subt : Cont → Cont → Set
  _,ₜ_     : Subt Δ Γ → Tm Δ → Subt Δ (Γ ▹ₜ)
_[_]ᵗ     : Tm Γ → Subt Δ Γ → Tm Δ
var vz [ γ ,ₜ a ]ᵗ = a
var (vs x) [ γ ,ₜ a ]ᵗ = var x [ γ ]ᵗ
fun ar ts [ γ ]ᵗ = fun ar (ts [ γ ]ᵗs)
data For  : Cont → Set
  _⊃_     : For Γ → For Γ → For Γ
  ∀       : For (Γ ▹ₜ) → For Γ
_[_]ᶠ     : For Γ → Subt Δ Γ → For Δ
(A ⊃ B) [ γ ]ᶠ = A [ γ ]ᶠ ⊃ B [ γ ]ᶠ
∀ A [ γ ]ᶠ = ∀ (A [ γ ⁺ ]ᶠ)
data Conp : Cont → Set
  _▹ₚ_    : Conp Γ → For Γ → Conp Γ
data Pf   : (Γ : Cont)(Γₚ : Conp Γ) → For Γ → Prop
data Subp : (Γ : Cont) → Conp Γ → Conp Γ → Prop
Con = Σ Cont Conp
Sub = λ (Δ,Δₚ) (Γ,Γₚ) → Σ(γ : Subt Δ Γ).Subp Δ Δₚ (Γₚ [ γ ]Conp)
    \end{minted}
\end{frame}

\begin{frame}{Összegzés}
  \begin{itemize}
    \item Absztrakt módon formalizált elsőrendű logika.
    \item Kvóciensek nélküli szintaxis.
    \item Több szemantika.
  \end{itemize}
  \vspace{1cm}
  További munkák:
  \begin{itemize}
    \item További szemantikák: pl. Kripke.
    \item Iterátor implementációja és inicialitás bizonyítása.
  \end{itemize}
\end{frame}

\begin{frame}[fragile]{Bizonyítás példa}
  \vspace*{-1.5em}
  \begin{align*}
    \texttt{A ⊃ B ⊃ A}
  \end{align*}
  Informális bizonyítás:
  {\fontsize{8pt}{8pt}\selectfont
    \[
    \infer[\text{⊃in}]{\vdash A ⊃ (B ⊃ A)}{
    \infer[\text{⊃in}]{A \vdash B ⊃ A}{
        \infer[\text{\tiny W-L}]{A, B \vdash A}{
            A \vdash A \quad \text{\tiny [Ax]}
        }
    }
    }
    \]
  }

  Bizonyítás a szintaxisban:
  \begin{minted}{agda}
    A⊃B⊃A : {A B : For ◇} → Pf ◇ (A ⊃ B ⊃ A)
    A⊃B⊃A = ⊃in (⊃in (qₚ [ pₚ ]ᵖ))
  \end{minted}
  Bizonyítás bármely modellben:
  \begin{minted}{agda}
    A⊃B⊃A : {A B : For ◇} → Pf ◇ (A ⊃ B ⊃ A)
    A⊃B⊃A {A}{B} =
      ⊃in (
        substP 
          (λ x → Pf (◇ ▹ₚ A) (x))
          (⊃[] ⁻¹)
          (⊃in {◇ ▹ₚ A}{B [ pₚ ]ᶠ}{A [ pₚ ]ᶠ} (qₚ [ pₚ ]ᵖ))
      )
  \end{minted}
\end{frame}

\end{document}
